\section{Sistema de Detección Rápida}

El cambio de color del agua dado por las nanopartículas funcionalizadas tras la
adsorción del arsénico puede ser bien aprovechado, pues permite la medición
\emph{in-situ} de contaminantes en componentes acuosos.

La primera estrategia a implementar es tomar la teoría que nos ofrece el
artículo y llevarlo a un proceso comercial, es decir, vender sobres/empaques con
nanoparticulas funcionalizadas. La idea consta en que el cliente destapa el
sobre y vierte el contenido (\emph{in-situ}) sobre el agua que desee analizar,
permitiéndole realizar una medida cualitativa de la muestra de manera visual.

Es prudente realizar pruebas que verifiquen si la funcionalización es estable,
es decir, si las nanoparticulas perduran unidas a los aminoácidos con el tiempo.
Ademas, también es necesario conocer en que manera se ven afectadas las
conexiones por causa de radiación electromagnética, pues podrían haber efectos
adversos que obstaculicen el proceso. Para ello podrían llegarse a plantear el
uso de empaques más sofisticados que, por ejemplo, no permitiese la entrada de
luz.

La segunda estrategia busca obtener una medida mas cuantitativa: Buscando
alternativas a un  equipo de caracterización UV-Vis se puede proponer un equipo
que funcione como una cámara oscura. El sistema ya vendría incorporado con
nanoparticulas funcionalizadas adyacentes a un diodo de un lado y una fotocelda
del otro. Tras la agregación de agua podríamos proceder a la inspección de picos
en longitud de onda debidos a resonancia de plasmones localizados de superficie.

Así mismo, como equipo de físicos, buscaríamos alternativas que remplacen el uso
de plasmones localizados de superficie (nanoparticulas) por el uso de plasmones
propagantes de superficie (laminas). Tuvimos como idea usar estas laminas
metálicas como recubrimiento interno de un envase donde se añadiría agua
contaminada con arsénico.

Sin embargo surgen algunas complicaciones puesto que aun haría falta adicionar
nanoparticulas que generen los enlaces, esto se debe a que la agregación no
seria posible únicamente con el recubrimiento interno entre puntos diferentes de
la lamina. En ese caso, ya sea cambiando el metal, el aminoácido con el que
funcionalizamos, o el metal a adsorber, las opciones de innovación se reducen
drásticamente.
